% !TEX root = ../main.tex

\frontchapter{Abstract (mini)}

The bulk-synchronous parallel (BSP) paradigm---powering large-scale scientific simulations and
distributed ML training---poses unique efficiency challenges.
These workloads impose periodic global synchronization across 10,000s of ranks, which compounds
small inefficiencies into cluster-wide bottlenecks.
Many arise from conditions unknowable before runtime, such as load imbalance and performance
pathologies, and can only be addressed through adaptation.
Efficiency in production therefore hinges on \emph{adaptation latency}: how quickly conditions can
be observed, characterized, and responded to.
This thesis argues that general-purpose infrastructure
for adaptive feedback loops
is the missing piece in today's BSP systems.

We develop this claim through two studies that instantiate the same primitive---a feedback loop
over application state---with starkly different outcomes.
On the data path, CARP implements an online loop for range-query-optimized data ingestion: it
periodically materializes a view of global key distributions via in-situ reduction and uses it to
recompute load-balanced range partitions as distributions evolve, matching the query performance of
a full sort with up to $4.9\times$ faster ingestion.
On the compute path, we study placement in adaptive mesh refinement codes, where apparent load
imbalance is confounded by faulty hardware, execution variability, and tuning artifacts.
Here the loop comprised multiple time-consuming cycles of
bulk trace collection, offline analyses, reconfiguration, and re-execution.
Isolating true placement effects consumed months even though individual mitigations typically took
days.
With reliable telemetry established, we design CPLX, a tunable placement policy that achieves up
to $21.6\%$ runtime improvement.
CARP demonstrates that online feedback loops are effective when they exist.
The AMR study demonstrates the cost when they do not.

Because these effects are intrinsically unpredictable, production BSP systems need an always-on,
programmable loop for observing and responding to runtime behavior.
Coarse views surface deviations, and the loop enables both steering into narrower views for
localization and triggering online interventions where viable.
We call this paradigm steerable observability.
ORCA realizes this via a BSP-aware tree-based overlay that
materializes runtime views through in-situ reduction and provides a timestep-consistent control
plane for coordinated interventions.
On a real AMR code at 4096 ranks, ORCA collects detailed traces at 1--2\% overhead versus
conventional HPC tracers at 56--81\% and eliminates 98.5--99.999\% of telemetry at source for
evaluated queries.
In an evaluation of its closed-loop capabilities using an LLM agent, a performance anomaly that
took months to diagnose manually is localized in 27 minutes.
